En un experiment fotoelèctric, il·luminem una superfície metàl·lica amb una llum verda que té una longitud d'ona de 546,1~nm. Observem que el potencial de frenada és de 0,376~V (tensió per la qual desapareix el corrent).

\begin{apartat}{1,25}
Determineu la funció de treball (treball d'extracció) d'aquesta superfície metàl·lica. Calculeu el llindar de freqüència per a l'extracció de fotoelectrons d'aquest metall.
\end{apartat}

\begin{apartat}{1,25}
Si il·luminem la superfície anterior amb una llum groga de 587,5~nm, determineu l'energia dels fotons incidents. Calculeu el potencial de frenada amb aquesta nova font de llum.
\end{apartat}

\dades{%
  $m_\mathrm{e} = 9,11\times 10^{-31}\ \mathrm{kg}$.\\
  $|e| = 1,602\times 10^{-19}\ \mathrm{C}$.\\
  $1\ \mathrm{eV} = 1,602\times 10^{-19}\ \mathrm{J}$.\\
  $c = 3,00\times 10^{8}\ \mathrm{m\,s^{-1}}$.\\
  $h = 6,63\times 10^{-34}\ \mathrm{J\,s}$.}
